% !TEX root = ../Main.tex
\chapter{斜率 $k$}

斜率 $k$ 描述一条直线的\textbf{倾斜程度}。设直线与 $x$ 轴正方向的夹角为倾斜角 $\alpha$，约定 $\alpha\in[0^\circ,180^\circ)$，则
\[
k=\tan\alpha,
\qquad
\bc
\alpha\ne 90^\circ, & k\ \text{存在};\\
\alpha=90^\circ, & k\ \text{不存在}.
\ec
\]
直线的常用写法是 $y=kx+b$，其中 $k$ 是斜率、$b$ 是纵截距。下面看两直线\textbf{平行}与\textbf{垂直}各对应斜率的什么关系。

\section{平行}

\state{两直线平行}{
    \[
    \ell_1\parallel\ell_2
    \quad\Longleftrightarrow\quad
    k_1=k_2,\ \underline{b_1\ne b_2}
    \quad(\text{两直线不可重合}).
    \]
}

原理很直接：平行就是\textbf{倾斜角相等}（$\alpha_1=\alpha_2$ 且 $\alpha_1,\alpha_2\ne 90^\circ$），倾斜角相等正切就相等，即 $k_1=k_2$。这里要补上 $b_1\ne b_2$，否则两条线会重合成同一条。

\rmk{两直线的斜率\textbf{都不存在}时（都竖直），它们也平行。}

\section{垂直}

\state{两直线垂直}{
    \[
    \ell_1\perp\ell_2
    \quad\Longleftrightarrow\quad
    k_1\cdot k_2=-1.
    \]
}

推导：设 $\ell_1$ 的倾斜角为 $\alpha_1$。两线垂直时，$\ell_2$ 的倾斜角比 $\ell_1$ 大 $90^\circ$，即
\[
\alpha_2=90^\circ+\alpha_1.
\]
于是
\[
\tan\alpha_2=\tan(90^\circ+\alpha_1)=-\frac{1}{\tan\alpha_1},
\qquad\text{即}\qquad
k_2=-\frac{1}{k_1},
\]
两边相乘即得 $k_1k_2=-1$。

\begin{center}
\begin{tikzpicture}[>=Stealth,scale=1.0]
  \draw[->] (-2.6,0) -- (2.8,0);
  \draw[thick,blue!60!black] (-1.6,-0.9) -- (2.4,1.35);
  \draw[thick,blue!60!black] (1.1,-1.3) -- (-1.1,1.95);
  \draw (0.9,0) arc (0:29:0.9);
  \node at (1.25,0.18) {\footnotesize$\alpha_1$};
  \draw (-0.55,0) arc (180:124:0.55);
  \node at (-0.55,0.4) {\footnotesize$\alpha_2$};
  \draw (0.18,0.27) -- (0.45,0.07) -- (0.25,-0.2);
\end{tikzpicture}
\end{center}

\rmkb{
    \textbf{注意点}：当一条直线 $k=0$（水平），另一条直线的斜率\textbf{不存在}（竖直）时，两条直线也互相垂直。此时不能套 $k_1k_2=-1$（一个斜率不存在），要单独识别。
}
